
% ════════════════════════════════════════════
% 模型准备（5.3.2.模型准备.tex，三级子块）—— 写作提示词见本文件末尾【本章节写作提示词】。
% ════════════════════════════════════════════
\subsubsection{模型准备}

\textbf{1.数据预处理}

本问误差参数由题面直接给出，无需常规缺失填补。为便于误差传播建模，首先将三类误差源统一为方位角域的标准差：车辆累计行驶 $s$ 米后，水平坐标误差为独立零均值正态随机变量，标准差 $\sigma_{\mathrm{pos}}(s)=0.3+0.02s$ 米；云台零位偏差标准差为 0.8 度，云台控制误差标准差为 0.6 度；目标位置精确已知，水平距离固定为 $r=40$ m。建立车辆定位误差向方位角偏差的换算系数，作为后续方差合成与概率反解的基础。

\textbf{2.算法选择}

针对本问误差传播与最大行驶距离反解，需在解析误差传播、数值积分与蒙特卡洛仿真之间择优，从计算精度、求解效率与普适性三个维度对比其适用性，如表\ref{tab:algo3}所示。

\begin{longtable}{>{\centering\arraybackslash}p{0.20\textwidth}>{\centering\arraybackslash}p{0.22\textwidth}>{\centering\arraybackslash}p{0.22\textwidth}>{\centering\arraybackslash}p{0.24\textwidth}}
  \caption{问题三候选算法对比}
  \label{tab:algo3} \\
  \toprule
  算法 & 计算精度 & 求解效率 & 普适性 \\
  \midrule
  \endfirsthead
  \caption{问题三候选算法对比（续）} \\
  \toprule
  算法 & 计算精度 & 求解效率 & 普适性 \\
  \midrule
  \endhead
  \bottomrule
  \endfoot
  解析误差传播 & 高（小误差） & 高 & 依赖线性化 \\
  数值积分 & 高 & 中 & 中 \\
  蒙特卡洛仿真 & 随样本量提高 & 低 & 强 \\
\end{longtable}

如表\ref{tab:algo3}所示，解析误差传播在小误差条件下精度高且可给出显式反解公式，蒙特卡洛仿真适用于验证。综合以上分析，本问以解析误差传播为主、蒙特卡洛仿真验证为辅进行求解。
